Nous utilisons des vues dans le base de donnée, afin de simplifier l'accès aux données complexes.
Pour la taxonomie des genres, nous nous basons sur une vue récursive qui nous permet pour chaque genre, d'avoir l'ensemble des sous genre de celui-ci.
\begin{lstlisting}[language=SQL, title={specializes\_transitive\_closure}]
CREATE VIEW specializes_transitive_closure AS
    WITH RECURSIVE specializes_transitive(genre, subgenre) AS (
        SELECT genre, subgenre FROM specializes
        UNION ALL
        SELECT st.genre, s.subgenre
        FROM specializes_transitive st, specializes s
        WHERE st.subgenre = s.genre
    )
    SELECT * FROM specializes_transitive;
\end{lstlisting}

Pour les statistiques sur les CDs, nous calculons d'abord des statistiques avec des fonctions utilitaires simples, puis pour les genres,
nous faisons l'union entre les genres présents sur les chansons du CD et tous les sous-genres de ces genres la, en utilisant la vue précédente.
\begin{lstlisting}[language=SQL, title={cd\_statistics}]
DROP VIEW IF EXISTS cd_statistics;
CREATE VIEW cd_statistics AS
    SELECT
        cd.cd_number,
        cd_get_total_song_duration(cd.cd_number) AS total_duration,
        cd_get_min_song_duration(cd.cd_number) AS min_duration,
        cd_get_max_song_duration(cd.cd_number) AS max_duration,
        cd_get_average_song_duration(cd.cd_number) AS average_duration,
        cd_playlist_count(cd.cd_number) AS playlist_count,
        (
            SELECT GROUP_CONCAT(all_genres.genre ORDER BY all_genres.genre SEPARATOR ', ')
            FROM (
                SELECT song.genre
                FROM song
                WHERE song.cd_number = cd.cd_number
                UNION
                SELECT stc.genre
                FROM song
                JOIN specializes_transitive_closure stc ON stc.subgenre = song.genre
                WHERE song.cd_number = cd.cd_number
            ) AS all_genres
        ) AS genres
    FROM cd;
\end{lstlisting}

Pour le dashboard des CDs, nous commençons par générer toutes les dates entre le tout premier et le tout dernier événement. Ensuite pour chaque
élément du produit cartésien de ces dates et des CD, on calcule le nombre de fois que ce CD est utilisé ce jour la.
Notons que si il y a un grand écart entre les deux dates et que la base de donnée contient beaucoup de CDS,
le produit cartésien générera une très grande quantité de tuples.
\begin{lstlisting}[language=SQL, title={cd\_dashboard}]
CREATE VIEW cd_dashboard AS
    WITH RECURSIVE dates(date) AS (
        SELECT MIN(date) AS date FROM event
        UNION ALL
        SELECT DATE_ADD(date, INTERVAL 1 DAY)
        FROM dates
        WHERE date < (SELECT MAX(date) FROM event)
    )
    SELECT
        dates.date,
        cd.cd_number,
        cd.copies,
        (
            SELECT COUNT(DISTINCT event.id)
            FROM event
            JOIN contains ON event.playlist = contains.playlist
            WHERE contains.cd_number = cd.cd_number AND event.date = dates.date
        ) AS copies_used
    FROM dates
    CROSS JOIN cd;
\end{lstlisting}

